\documentclass[aps,pra,twocolumn,superscriptaddress,longbibliography]{revtex4-2}

\usepackage{amsmath,amssymb,amsthm,graphicx,xcolor,bm,braket,dcolumn,booktabs}
\usepackage{hyperref}

\newtheorem{observation}{Observation}
\newcommand{\gvq}[2]{#2}
\newcommand{\gvc}[2]{#2}
\newcommand{\mHa}{\,\mathrm{mHa}}

\begin{document}

\title{Water at Twenty-Four Functions:\\
Slater-Type Bases Set on the Free Atom Against cc-pVDZ,\\
With the Qubit Cost}
\thanks{Paper 62 in the Geometric Vacuum series}

\author{J.~Loutey}
\affiliation{Independent Researcher, Kent, Washington}

\date{October 31, 2026}

\begin{abstract}
\emph{Central result.}  At a fixed basis size of twenty-four functions,
the size of cc-pVDZ on water, we ask how much of the Gaussian basis's
correlated energy a one-centre Slater-type design recovers when its
exponents are set on the free atoms, and what that costs in qubit
resources.  \textbf{[MEASURED]} A Slater double-zeta-plus-polarization
basis built by cc-pVDZ's own design rule sits
$\gvq{p62_dzp_vs_ccpvdz_mha}{40.7}\mHa$ below cc-pVDZ at frozen-core
CCSD(T), of which $\gvq{p62_dzp_vs_ccpvdz_scf_mha}{28.8}\mHa$ is
self-consistent field.  \textbf{[MEASURED, negative]} Every exponential
kit with fixed or shared exponents loses: the hydrogenic $Z/n$ rule by
$\gvq{p62_hyd_above_ccpvdz_mha}{196}\mHa$, a single Coulomb--Sturmian
scale per atom by more than one hartree, and a two-scale Sturmian kit
by $\gvq{p62_two_scale_above_ccpvdz_mha}{130}\mHa$; the hydrogenic
loss is the fixed radial node structure, not the decay rate.
\textbf{[MEASURED, pre-registered]} Releasing the oxygen exponents one
at a time, every value set by the free oxygen atom's radial Hartree--Fock
energy with nothing tuned on water, closes that gap: with five oxygen
exponents the water energy lands $\gvq{p62_r5_below_ccpvdz_mha}{3.1}\mHa$
below cc-pVDZ, and $\gvq{p62_r5h_below_ccpvdz_mha}{7.0}\mHa$ below with
hydrogen exponents set on H$_2$.  The split of the $2p$ from the $3p$
exponent carries $\gvq{p62_sl1_ratio}{0.900}$ of the water gain, as it
carries $0.902$ of the atom's; the water gain is
$\gvq{p62_sl2_ratio}{0.516}$ of the atom's.  Both were predicted before
the water energies were computed.  The one non-monotone rung is
diagnosed: its correlation loss transfers from the atom one to one and is
mostly a $3d$ exponent the atomic objective never sees, while its
Hartree--Fock gain is two opposite pieces that cancel on water.
\textbf{[MEASURED]} The Jordan--Wigner Pauli count is the same for every
twenty-four-function basis, $\gvq{h2o_ccpvdz_pauli}{107\,382}$ terms on
46 frozen-core qubits, because it is fixed by the orbital count and the
$C_{2v}$ pattern; the 1-norm rises with accuracy, from 405 to 706~Ha
across the kits.  \emph{Caps.}  One geometry, one method, one molecule.
The water-tuned Slater basis stays $\gvq{p62_dzp_below_r5h_mha}{34}\mHa$
below the best atom-set rung.  \textbf{[MEASURED, pre-registered]} The
matched-accuracy comparison brackets the Gaussian cost from both sides:
aug-cc-pVDZ at 41 functions stays $\gvq{p62_aug_above_dzp_mha}{8.0}\mHa$
short of the Slater DZP while already costing
$\gvq{p62_aug_pauli_ratio}{9.2}\times$ the Pauli terms and
$\gvq{p62_aug_lambda_ratio}{2.7}\times$ the 1-norm on 80 qubits against
46, and the crossing lies between it and cc-pVTZ.  Nothing here is a
GeoVac lattice structure: the
bases are tuned exponential functions evaluated through Gaussian fits,
and the angular sparsity of the one-centre framework does not enter.
\end{abstract}

\maketitle

\section{Introduction}
\label{sec:intro}

The atomic sector of the Geometric Vacuum programme is built on
exponential radial functions: hydrogenic orbitals with decay $Z/n$, and
Coulomb--Sturmians with one shared scale per centre
\cite{avery2006,avery2017}.  Its molecular work has either kept those
functions and paid for it in accuracy (the composed architecture,
Paper~17), or left them for prolate-spheroidal and explicitly correlated
constructions on two centres (Papers~11, 12, 19).  Water is the first
molecule in the series where the integrals over exponential functions
on three centres are evaluated to working precision (Paper~59 and the
shared-centre channels of CHANGELOG v5.20.10), and so the first where
the question of this paper can be asked cleanly.

The question is narrow on purpose.  Fix the basis size at twenty-four
functions, which is what cc-pVDZ \cite{dunning1989} spends on water.
Fix the geometry and the correlation method.  Then ask three things.
First, does a Slater-type basis built by cc-pVDZ's own design rule beat
cc-pVDZ at that size, and by how much at the mean-field and the
correlated level separately?  Second, do the exponential functions the
one-centre framework actually uses, with their fixed or shared
exponents, survive the same test?  Third, if they do not, how many
distinct radial scales does oxygen need, and can those scales be set on
the free atom rather than on the molecule?

The third question is the one with content.  A basis whose exponents
are tuned on the molecule being computed is a fit.  A basis whose
exponents are set on the free atom and then carried to the molecule
tests a transfer.  The scale ladder of Sec.~\ref{sec:ladder} was
pre-registered as a transfer test: the predictions were written and
committed before any water energy of the new rungs was computed, and
the record is reproduced in Sec.~\ref{sec:ladder} with its disposition.

Two things this paper is not.  It is not a GeoVac lattice result.  The
bases here are tuned Slater-type functions evaluated through Gaussian
fits, and the angular (Gaunt) sparsity of the one-centre framework
\cite{paper22} plays no part in any number below; Sec.~\ref{sec:qubit}
says what the qubit count does and does not depend on.  And it is not a
claim about the basis-set limit.  Everything is measured at one size,
one geometry and one method, and the caps are listed in
Sec.~\ref{sec:limits}.

\section{Solver and validation}
\label{sec:solver}

\paragraph{Integrals.}  All one- and two-electron integrals are
evaluated by the McMurchie--Davidson scheme \cite{mcmurchie1978} over
contracted Cartesian Gaussians, through the project's
\texttt{noci\_engine} module.  Gaussian bases are used as published.
Slater-type primitives $r^{n-1}e^{-\zeta r}Y_{lm}$ enter as Gaussian
fits: one shape is fitted per $(l,n)$ and number of Gaussians, by
least squares on the radial function, and then scaled by $\zeta^{2}$ to
the required exponent \cite{stewart1970}.  Because the shape is fitted
once and scaled, the fit error of a given $(l,n)$ is the same at every
exponent; this matters for the ladder of Sec.~\ref{sec:ladder}, where
rungs differ only in exponents.  Twelve Gaussians per valence primitive
and eight per polarization function are used throughout.  The fitted
basis is a legitimate Gaussian basis with exact integrals, so every
Hartree--Fock energy below is a strict upper bound for that fitted
basis; the fit's convergence toward the true Slater function is checked
in Sec.~\ref{sec:dzp}.

\paragraph{Correlation.}  The correlated solver is the project's own:
restricted Hartree--Fock, spin-orbital coupled-cluster singles and
doubles, and the perturbative triples correction
\cite{purvis1982,raghavachari1989}, with an optional frozen oxygen $1s$.
It is in \texttt{tests/test\_water\_ccsd.py}.  (An external package was
not used because none builds on the host; the solver's independence is
also what makes the checks below meaningful.)  It is validated where
its answer is known:
\begin{itemize}
\item it reproduces the published STO-3G water reference of the
Crawford programming projects \cite{crawford_projects} to
$2\times10^{-12}$~Ha in the triples correction and
$4\times10^{-10}$~Ha in the CCSD correlation energy;
\item CCSD equals full configuration interaction for two electrons, to
$10^{-9}$~Ha;
\item frozen-core CCSD equals active-space full CI on a case where both
are computable;
\item nine planted wrong answers (sign flips, dropped diagrams, a wrong
denominator) were each caught by the tests.
\end{itemize}

\paragraph{Atomic reference.}  Oxygen exponents are set on the free
atom by an exact one-centre radial average-of-configuration
Hartree--Fock over Slater primitives, with the Slater $F^{k}$ and
$G^{k}$ integrals in closed form
\cite{slater1960}.  The code gives helium's single-zeta energy
$-729/256$~Ha exactly, and for neon its energy is the limit that the
three-dimensional Gaussian-fit route approaches from above as the
number of Gaussians grows ($-4.9$, $-0.20$, $-0.09$~mHa at 8, 12, 16
Gaussians).  An earlier open-shell route through fractional-occupation
SCF orbitals sat $10.7$~mHa above the variational average energy, because
its orbitals are not variational for that energy; it was replaced by the
radial code and is not used.

\paragraph{Geometry and conventions.}  $R_{\rm OH}=1.809$~bohr,
$\angle\mathrm{HOH}=104.5^{\circ}$.  Energies are in hartree.  ``fc''
means the oxygen $1s$ is frozen in the correlation treatment; cc-pVDZ is
a valence-correlation basis, so the frozen-core energy is the fair
yardstick, and it is the scored quantity throughout.  All-electron
values are reported alongside where they were computed.

\section{Equal size, equal design rule: Slater DZP against cc-pVDZ}
\label{sec:dzp}

The fairness rule is cc-pVDZ's own design, with Slater primitives in
place of Gaussian ones.  cc-pVDZ on oxygen is a general contraction of
atom-optimized primitives to $[3s2p1d]$, with the contracted functions
being the atomic $1s$ and $2s$ orbitals plus the most diffuse $s$
primitive, the atomic $2p$ plus the most diffuse $p$ primitive, and one
$d$ polarization function.  The Slater basis here does the same:
\begin{itemize}
\item O: four $s$ primitives ($1s,1s',2s,2s'$) and three $p$ primitives
optimized on the atom's average-of-configuration Hartree--Fock energy,
contracted to $\{$atomic $1s$, atomic $2s$, most diffuse $s\}$ and
$\{$atomic $2p$, most diffuse $p\}$.  The optimized exponents are
$9.317, 6.765 \mid 2.710, 1.694$ for $s$ and $5.719, 2.846, 1.450$ for
$p$, with atom energy $-74.768205$~Ha.
\item H: $1s,1s'$ optimized on H$_2$ RHF at $R=1.4$~bohr: $1.212$,
$0.627$.
\item Polarization (O $3d$, H $2p$): one exponent each, optimized on
frozen-core CCSD of water, for both bases alike.  The cc-pVDZ $d$ and
$p$ exponents re-optimized the same way move from $1.185$ to $1.312$ and
from $0.727$ to $0.605$; the Slater optima are O $3d$ $2.756$ and H
$2p$ $1.579$.  (A Slater $3d$ is diffuse at a given $\zeta$, with
$\langle r\rangle=7/2\zeta$, so these look large next to Gaussian
values.)
\end{itemize}
Polarization is the one place where water enters the design, and it
enters both bases identically.  A negative result fixed the contraction
rule: an uncontracted $3s2p$ set optimized on the atom spends its
second $s$ function on the core ($10.03$, $8.33$, $2.32$), leaving the
valence $s$ single-zeta.  That is why cc-pVDZ uses general contraction,
and why the Slater basis does too.

\begin{table}[t]
\caption{Water at 24 functions, $R_{\rm OH}=1.809$~bohr,
$104.5^{\circ}$.  fc = oxygen $1s$ frozen in the correlation treatment.
The cc-pVTZ row (58 functions) is the anchor for the Gaussian ladder.
Energies in hartree.}
\label{tab:dzp}
\begin{ruledtabular}
\begin{tabular}{lrrr}
Basis & \multicolumn{1}{c}{RHF} & \multicolumn{1}{c}{CCSD(T) fc} &
\multicolumn{1}{c}{CCSD(T) ae} \\
\hline
cc-pVDZ (standard) & -76.026795 & \gvq{p62_ccpvdz_fc}{-76.241025} & -76.243143 \\
cc-pVDZ, pol.\ re-opt.\ & -76.027257 & -76.242201 & -76.244383 \\
Slater DZP (12 G) & \gvq{p62_dzp_rhf}{-76.055573} & \gvq{p62_dzp_fc}{-76.281739} & -76.285720 \\
Slater DZP (16 G) & -76.055644 & -76.281811 & -76.285792 \\
cc-pVTZ (58 fns) & \gvq{p62_ccpvtz_rhf}{-76.057162} & \gvq{p62_ccpvtz_fc}{-76.332191} & \\
\end{tabular}
\end{ruledtabular}
\end{table}

\textbf{[MEASURED]} Table~\ref{tab:dzp} gives the result.  At
frozen-core CCSD(T) the Slater DZP basis is
$\gvq{p62_dzp_vs_ccpvdz_mha}{40.7}\mHa$ below cc-pVDZ at equal size,
and $39.5$~mHa below cc-pVDZ with its polarization re-optimized on the
same footing.  Of the gain, $\gvq{p62_dzp_vs_ccpvdz_scf_mha}{28.8}\mHa$
is at the Hartree--Fock level and $12$~mHa is additional correlation
energy.  The evaluator is converged: going from twelve to sixteen
Gaussians per function moves every energy by $0.07$~mHa.

\paragraph{Place on the Gaussian ladder.}  The cc-pVTZ anchor makes the
split readable.  The Slater DZP RHF energy is $1.6$~mHa above
cc-pVTZ's, so it covers about $94\%$ of the DZ-to-TZ mean-field gap.
Its frozen-core CCSD(T) energy is $50.5$~mHa above cc-pVTZ's, so it
covers about $45\%$ of the DZ-to-TZ correlated gap.  The reading is the
expected one: the exponential shape buys the mean-field part, the cusp
at the nucleus and the exponential tail, while correlation energy is
bought by the number of correlating functions, which is the same in
both bases.  That the mean-field advantage of Slater functions exists
is textbook; what is measured here is its size at equal count under an
identical design rule, and the fact that the correlated advantage is
less than a third of it.

\section{The framework's own exponential kits, negative}
\label{sec:native}

The $41$~mHa of Sec.~\ref{sec:dzp} could come from the exponential
shape or from the exponent tuning.  The one-centre framework's radial
rules fix the exponents, so they separate the two.  Two rules and two
remedies were tested, at the same twenty-four functions (shells
$n\le3$ on O, $n\le2$ on H), the same geometry and the same solver, with
no exponent or polarization tuned on water in any of them.

\begin{description}
\item[Hydrogenic $Z/n$, bare $Z$.]  Every function is a hydrogen-like
orbital of the nuclear charge, decay $Z/n$, with its fixed radial
polynomial.  This is the production rule of the composed builder
(Paper~17) and the two-centre engine (Paper~58).  Zero parameters.  A
slow test checks that each function is a one-electron eigenfunction at
$-Z^{2}/2n^{2}$ and that same-$l$ functions are orthogonal; it caught
two planted coefficient errors.
\item[Hydrogenic with the composed screening.]  O $1s$ at $Z=8$; O
shells $n\ge2$ at $Z_{\rm eff}=6$, which is $Z$ minus the $1s^{2}$ core
as the composed classifier assigns it.  Zero parameters.  This maps
the composed core and valence blocks onto one function set; the
composed builder itself starts its valence block at $n=1$, so the
mapping is ours.
\item[Coulomb--Sturmian, one scale per atom.]  Avery's basis
\cite{avery2006}: one shared exponent $k$ on each centre.  $k_{\rm H}=1$
makes the H $1s$ exact.  $k_{\rm O}=3.781$ minimizes the radial oxygen
energy, and the best that one shared scale can do for the atom is
$-71.406$~Ha, $3.4$~Ha above the Slater primitives' $-74.768$.
\item[Coulomb--Sturmian, two scales.]  A core $1s$ at $k_{c}$ plus
Sturmians sharing $k_{v}$ ($s$ with $n\le2$, $p$ with $n\le3$, one
$3d$) on oxygen, fourteen functions; $(k_{c},k_{v})=(7.793, 2.281)$
minimize the radial oxygen energy ($-74.5027$~Ha); H at $k=1$.  This is
a simplified relative of Avery's generalized Sturmians, not his
isoenergetic construction (Paper~60).
\end{description}

\begin{table}[t]
\caption{The exponential kits at 24 functions, frozen-core CCSD(T),
same geometry and solver as Table~\ref{tab:dzp}.  STO-3G (7 functions)
for reference.  Every 24-function kit has
$\gvq{h2o_ccpvdz_pauli}{107\,382}$ Jordan--Wigner Pauli terms on 46
qubits (Sec.~\ref{sec:qubit}); the frozen-core 1-norm is in hartree.}
\label{tab:native}
\begin{ruledtabular}
\begin{tabular}{lrrrr}
Kit & Fns & \multicolumn{1}{c}{RHF} & \multicolumn{1}{c}{CCSD(T) fc} &
\multicolumn{1}{c}{$\lambda$ fc} \\
\hline
STO-3G & 7 & -74.9629 & -75.0123 & \\
Sturmian, one $k$ per atom & 24 & -74.2879 & -74.5603 & 848.8 \\
hydrogenic, $Z_{\rm eff}=6$ ($n\ge2$) & 24 & -75.7466 & \gvq{p62_hyd_zeff_fc}{-75.9471} & 405.1 \\
hydrogenic, bare $Z$ & 24 & -75.8397 & \gvq{p62_hyd_fc}{-76.0452} & 510.8 \\
Sturmian, two scales & 24 & -75.8903 & \gvq{p62_two_scale_fc}{-76.1113} & 620.0 \\
cc-pVDZ & 24 & -76.0268 & \gvq{p62_ccpvdz_fc}{-76.2410} & 674.3 \\
Slater DZP & 24 & -76.0556 & \gvq{p62_dzp_fc}{-76.2817} & 706.3 \\
\end{tabular}
\end{ruledtabular}
\end{table}

\textbf{[MEASURED, negative]} Table~\ref{tab:native} gives the result.
No kit with fixed or shared exponents reaches cc-pVDZ at twenty-four
functions.  The bare-$Z$ hydrogenic kit is
$\gvq{p62_hyd_above_ccpvdz_mha}{196}\mHa$ above it and $236$~mHa above
the Slater DZP.  The one-scale Sturmian kit, at twenty-four functions,
is worse than STO-3G at seven.  So the $41$~mHa of Sec.~\ref{sec:dzp}
came from exponent choice, not from the hydrogen-like shape.

\paragraph{Mechanism, hydrogenic.}  Bare $Z=8$ puts the oxygen $n=2$
shell at $\zeta=4$, about twice as compact as the tuned valence
($1.45$ to $2.85$).  The valence then has to live in the $n=3$
functions ($\zeta=2.67$), which leaves almost nothing to correlate with.
The core is fine ($\zeta=8$ against $7.7$ to $9.3$ tuned).  Screening
is the obvious remedy and it makes things worse, by $98$~mHa.  The
diagnosis comes from the oxygen atom in the kits as built: bare
$-74.557$, screened $-74.420$, and the same two kits with their
polynomial contractions released give $-74.757$, an upper bound on what
the contracted set could reach.  So between $200$ and $340$~mHa is lost
to the fixed hydrogenic node structure.  Each hydrogenic function
carries a radial polynomial whose nodes are set for orthogonality to the
$1s$ of the \emph{same} charge.  Bare $Z$ keeps every shell mutually
consistent; $Z_{\rm eff}=6$ puts the valence nodes in the wrong place
relative to the $Z=8$ core.  The contraction, not the decay rate, is
the binding constraint.

\paragraph{Mechanism, Sturmian.}  One scale cannot serve a core near
$\zeta=8$ and a valence near $\zeta=2$ at once.  Sturmians pay no
contraction penalty, since a shared-$k$ set spans the same space as its
primitives, and the two-scale kit is accordingly the best of the four,
at $\gvq{p62_two_scale_above_ccpvdz_mha}{130}\mHa$ above cc-pVDZ.
The remedy the Sturmian literature names is several scales
\cite{avery2006}; whether two is the cap or merely the first rung is
what Sec.~\ref{sec:ladder} tests.

\paragraph{An observation, not a claim.}  Across the twenty-four
function kits other than the one-scale Sturmian, the frozen-core
1-norm rises with accuracy: $405$, $511$, $620$, $674$, $706$~Ha.  A
low 1-norm at equal size here tracks a worse description, not a cheaper
one.  Sec.~\ref{sec:qubit} returns to this.

\section{The scale ladder: how many oxygen exponents, set on the atom}
\label{sec:ladder}

\subsection{Design, fixed before any water rung was run}

Twenty-four functions throughout.  O: $s$ primitives ($1s,1s,2s$), $p$
primitives ($2p,3p$), and one $3d$ tied to the $3p$ exponent.  H:
$1s,2s,2p$.  Rungs differ only in which oxygen primitives share an
exponent (Table~\ref{tab:ladder_design}).  Every oxygen exponent set
minimizes the radial average-of-configuration Hartree--Fock energy of
the free oxygen atom (Sec.~\ref{sec:solver}), Nelder--Mead from 41
starts.  Hydrogen exponents are $k=1$, except on the last rung, where
$1s,2s,2p$ are each set by H$_2$ RHF at $R=1.4$~bohr.  Nothing is tuned
on water.  So any water gain is transfer, not fitting.  The $R2$ rung
is the two-scale Sturmian kit of Sec.~\ref{sec:native}; rebuilt through
the ladder's driver it reproduces that kit's energies to $10^{-4}$~Ha
(the pre-registered control SL-0).

\begin{table}[t]
\caption{The ladder.  Letters mark which oxygen primitives share an
exponent.  Atom energies are the radial average-of-configuration
Hartree--Fock minima that set the exponents; the seven-primitive Slater
atom of Sec.~\ref{sec:dzp} is $-74.76820$~Ha, $8.2$~mHa below $R5$.}
\label{tab:ladder_design}
\begin{ruledtabular}
\begin{tabular}{llcrr}
Rung & O groups $(1s,1s,2s,2p,3p)$ & $N_\zeta$ & \multicolumn{1}{c}{Atom}
& \multicolumn{1}{c}{Step (mHa)} \\
\hline
$R2$ & $a,b,b,b,b$ & 2 & -74.50266 & \\
$R3$ & $a,b,b,c,c$ & 3 & -74.50605 & 3.4 \\
$R4$ & $a,b,c,d,d$ & 4 & -74.52793 & 21.9 \\
$R5$ & $a,b,c,d,e$ & 5 & -74.75997 & 232.0 \\
$R5H$ & as $R5$; H from H$_2$ & 5 & & \\
\end{tabular}
\end{ruledtabular}
\end{table}

The atom ladder is the control that frames the predictions.  Its total
gain $R2\to R5$ is $257.3$~mHa, and the $p$ split ($R4\to R5$), which
releases the $2p$ exponent from the $3p$, carries $0.902$ of it.  Two
transfer ratios were measured on cases the ladder does not score: the
clean step from one shared oxygen scale to two (the two Sturmian kits of
Sec.~\ref{sec:native}) transfers at $0.50$, and the confounded step from
the two-scale kit to the Slater DZP at $0.64$, confounded upward because
the DZP also has two hydrogen $s$ exponents and water-tuned
polarization.

\subsection{Predictions, as registered}

The record \cite{sl_record} was committed on its own before any water
energy of $R3$, $R4$, $R5$ or $R5H$ was computed.  With $W(R)$ the
frozen-core CCSD(T) water energy of rung $R$ and $G=W(R2)-W(R5)$:
\begin{description}
\item[SL1 (gated).]  $(W(R4)-W(R5))/G\ge0.6$: the $p$ split dominates
on water too.  Stated probability $0.75$.
\item[SL2 (gated).]  $0.3\le G/257.3\mHa\le0.8$, i.e.\ $W(R5)$ between
$-76.1885$ and $-76.3172$~Ha.  The band is wide on purpose, because the
clean control measured a core step and the ladder's main step is
valence.  Probability $0.7$.
\item[SL3 (reported).]  $W(R5)$ stays above the Slater DZP.  Direction
only.  Probability $0.6$.
\item[SL4 (reported).]  $W(R5H)<W(R5)$.  Direction only.  Probability
$0.7$.
\end{description}
Gate: GO if SL1 and SL2 both hold.  Also reported, not predicted:
$W(R5)$ against cc-pVDZ, which at the clean transfer ratio lands on top
of it ($-76.240$ against $-76.241$) and so was declared a coin flip in
advance.

\subsection{Result}

\begin{table}[t]
\caption{Water on the ladder, 24 functions, same geometry and solver.
Steps are rung-to-rung changes, positive downward.  The atom step is the
radial HF gain of Table~\ref{tab:ladder_design}.}
\label{tab:ladder}
\begin{ruledtabular}
\begin{tabular}{lrrrrr}
Rung & \multicolumn{1}{c}{RHF} & \multicolumn{1}{c}{CCSD(T) fc} &
\multicolumn{1}{c}{CCSD(T) ae} & \multicolumn{1}{c}{Water step} &
\multicolumn{1}{c}{Atom step} \\
\hline
$R2$ & -75.890293 & \gvq{p62_two_scale_fc}{-76.111337} & -76.120634 & & \\
$R3$ & -75.890159 & \gvq{p62_r3_fc}{-76.108484} & -76.117764 & -2.85 & 3.4 \\
$R4$ & -75.905860 & -76.124580 & -76.137966 & 16.10 & 21.9 \\
$R5$ & -76.036106 & \gvq{p62_r5_fc}{-76.244112} & -76.258281 & 119.53 & 232.0 \\
$R5H$ & -76.045000 & \gvq{p62_r5h_fc}{-76.248015} & -76.262063 & 3.90 & \\
\end{tabular}
\end{ruledtabular}
\end{table}

\textbf{[MEASURED, pre-registered; GO]}  Table~\ref{tab:ladder}.
$G=132.78$~mHa.
\begin{itemize}
\item SL1 held: $(W(R4)-W(R5))/G=\gvq{p62_sl1_ratio}{0.900}$, against
the bound $0.6$ and the atom's $0.902$.
\item SL2 held: $G/257.3\mHa=\gvq{p62_sl2_ratio}{0.516}$, inside the
band; the clean control was $0.50$.
\item SL3 held in direction: $W(R5)$ is $37.7$~mHa above the Slater
DZP.  The record's estimate from the confounded ratio had been about
$6$~mHa; that estimate was well off, and the record says so.
\item SL4 held: $W(R5H)$ is $3.90$~mHa below $W(R5)$, against an H$_2$
gain of $4.06$~mHa, a transfer near $0.96$ on the hydrogen side.
\end{itemize}
Reported, not predicted: against cc-pVDZ at the same twenty-four
functions, $R5$ is $\gvq{p62_r5_below_ccpvdz_mha}{3.1}\mHa$ below and
$R5H$ is $\gvq{p62_r5h_below_ccpvdz_mha}{7.0}\mHa$ below.  The two-scale
kit had been $130$~mHa above.  So the gap of Sec.~\ref{sec:native} is
closed by releasing exponents set on the free atom and on H$_2$, with
nothing tuned on water.  RHF carries most of it: $145.8$~mHa from $R2$
to $R5$ at the mean-field level against $132.8$ at CCSD(T).

What this does and does not say.  $R5$ is a Slater-type basis with five
atom-set oxygen exponents, a $3d$ tied to the $3p$, and one hydrogen
exponent.  It is a better-chosen member of the same primitive family as
the two-scale kit, not a new structure.  The Slater DZP, with
water-tuned polarization and two hydrogen $s$ exponents, stays
$\gvq{p62_dzp_below_r5h_mha}{34}\mHa$ below $R5H$.  That $34$~mHa is
the price of setting everything on the atom.  The lesson of
Sec.~\ref{sec:native} stands: the case rests on one-centre structure,
not on better bricks; what moved is that the two-exponent cap, not the
Sturmian family, had been the $130$~mHa.

\subsection{The non-monotone rung, diagnosed}
\label{sec:r3}

$R3$ splits the oxygen $p$ exponent from the $s$-valence exponent.  The
atom gains $3.4$~mHa; water loses $2.85$~mHa.  The rung-by-rung ratios
of water step to atom step, $-0.84$, $0.735$, $0.515$, are therefore not
one transfer number per ladder.  A diagnostic, with its hypotheses
written before any new number was computed, located the cause
\cite{sl_record}.  Nothing was tuned on water: every substituted exponent
below is $R2$'s $2.28124$ or $R3$'s own value.

The water loss is correlation.  From $R2$ to $R3$, RHF moves $+0.13$~mHa,
the CCSD correlation $+2.57$ and the triples $+0.15$.  The one-centre
probe is the neutral oxygen atom in the same Gaussian-fitted kits, with
an unrestricted Hartree--Fock reference for the $M_S=1$ triplet
determinant and the same frozen-core CCSD(T).  On the atom the step
gains $2.84$~mHa at UHF and loses $2.72$~mHa of correlation, a wash of
$-0.12$~mHa.  Water loses the same $2.72$~mHa of correlation.  So the
correlation half of the step transfers one to one, and the Hartree--Fock
half does not transfer at all.

Where the correlation goes.  The ladder's objective is Hartree--Fock on
the atom, with no $d$ function.  It never sees the $3d$, and it sees the
$3p$ only through the occupied $2p$.  At $R3$ the $p$ family drifts
diffuse to suit the occupied $2p$, and the $3p$ and $3d$, which carry
most of the correlation at this size and want to be compact, ride along.
Putting the $3d$ alone back at $R2$'s exponent, with everything else at
$R3$, recovers $2.13$ of the $2.85$~mHa on water; the atom had predicted
about $2.1$ before that run.  About $80\%$ of the loss is the $3d$, a
passenger of the $3p$ by the ladder's design, with the rest split
between the $3p$ and the $s$ shift.

Why the Hartree--Fock gain cancels.  The atom gains at HF from both
moves of the step, the $p$ family going diffuse and the $s$ valence
tightening.  On water, the first costs $5.2$~mHa at RHF and the second
gains $5.1$~mHa, so the atom's gain cancels to $0.13$~mHa.  Water's RHF
is two to six times more sensitive than the atom's to the $3p$ and
$s$-valence exponents, and of the opposite sign for the $p$ family: the
free atom's $2p^{4}$ wants a more diffuse $p$ than the oxygen in water's
bonds.  The prediction's clause that RHF would stay within $0.2$~mHa of
$R3$ under each substitution failed for the $3p$ and $s$ substitutions,
by $6.7$ and $5.1$~mHa; the correlation clause held in all three.  That
failure is the mechanism.

The upshot for the ladder: $R3$'s ratio of $-0.84$ is a ratio of two
cancelling HF pieces and a transferring correlation loss, not a transfer
number, and the $R4$ and $R5$ rungs, which are HF-driven radial splits,
are untouched.  The ladder's one design flaw is identified: a $3d$ tied
to an exponent the atomic objective cannot see.  Setting it by an atomic
correlated criterion instead would change the ladder's rule and was not
done.

\section{The qubit cost}
\label{sec:qubit}

The second-quantized Hamiltonian of each basis was mapped to qubits by
the Jordan--Wigner transformation \cite{jordan1928} over RHF canonical
orbitals, with interleaved spin, terms kept at $|c|>10^{-10}$, the
identity counted in the Pauli number and excluded from the 1-norm.  The
counter is closed-form per string class and is validated against a full
Jordan--Wigner construction in OpenFermion on STO-3G water, where five
planted errors were each caught; its frozen-core cc-pVDZ count matches
the published count of Trenev \emph{et al.}\ \cite{trenev2025} exactly.

\begin{table}[t]
\caption{Qubit resources, Jordan--Wigner.  $\lambda$ is the Pauli
1-norm in hartree, identity excluded.  Frozen core (fc) drops the oxygen
$1s$.}
\label{tab:qubit}
\begin{ruledtabular}
\begin{tabular}{lrrrrrrr}
Basis & $Q$ fc & $N_{\rm Pauli}$ fc & \multicolumn{1}{c}{$\lambda$ fc} &
$Q$ ae & $N_{\rm Pauli}$ ae & \multicolumn{1}{c}{$\lambda$ ae} &
\multicolumn{1}{c}{CCSD(T) fc} \\
\hline
Slater DZP & 46 & \gvq{h2o_ccpvdz_pauli}{107\,382} & \gvq{p62_dzp_lambda_fc}{706.3} & 48 & 128\,793 & 754.6 & -76.2817 \\
cc-pVDZ & 46 & \gvq{h2o_ccpvdz_pauli}{107\,382} & \gvq{p62_ccpvdz_lambda_fc}{674.3} & 48 & 128\,793 & 717.0 & -76.2410 \\
cc-pVTZ & 114 & 3\,983\,652 & 7\,653.5 & 116 & 4\,281\,643 & 7\,858.6 & -76.3322 \\
\end{tabular}
\end{ruledtabular}
\end{table}

\textbf{[MEASURED]}  Table~\ref{tab:qubit}.  Three facts.
\begin{enumerate}
\item The Pauli count is identical in every twenty-four-function basis,
tuned or not, Gaussian or exponential.  It is fixed by the orbital count
and the $C_{2v}$ pattern of the canonical orbitals, which every basis
here shares, and not by the radial shape.  Radial shape can move an
energy by a hartree (Table~\ref{tab:native}) without moving a single
term.
\item The 1-norm is $\gvq{p62_lambda_ratio_pct}{4.7}\%$ higher for the
Slater DZP than for cc-pVDZ.  So the $41$~mHa of Sec.~\ref{sec:dzp}
costs about five percent in 1-norm and nothing in qubits or terms.
\item cc-pVTZ, the next Gaussian rung, costs $37\times$ the terms and
$11\times$ the 1-norm for the remaining $55\%$ of the correlated gap.
\end{enumerate}
The second fact is an equal-size statement, not a matched-accuracy
one.  The Slater DZP sits between cc-pVDZ and cc-pVTZ in energy, so the
matched-accuracy question is which Gaussian basis first reaches its
energy and what that basis costs.

\subsection{Matched accuracy}
\label{sec:matched}

Two intermediate Gaussian rungs were added to bracket the crossing from
both sides, and the outcome was pre-registered before either was run
\cite{ma_record}: cc-pVDZ's $s$ and $p$ shells with cc-pVTZ's
polarization shells (O $2d$, H $2p$; 35 functions; our construction, there
to bracket, not to recommend), and the standard aug-cc-pVDZ (41
functions).  The gated predictions were that neither reaches the Slater
DZP; the control bounds were measured on the neon atom first, where the
second polarization set is worth $26$~mHa, all correlation, and
augmentation $30$~mHa, a quarter of it mean-field.

\begin{table}[t]
\caption{Matched accuracy.  Same geometry, solver and counter as
Table~\ref{tab:qubit}.  ``vs DZP'' is the frozen-core CCSD(T) energy
relative to the Slater DZP's, positive above.}
\label{tab:matched}
\begin{ruledtabular}
\begin{tabular}{lrrrrr}
Basis & Fns & CCSD(T) fc & vs DZP (mHa) & $N_{\rm Pauli}$ fc ($Q$) & $\lambda$ fc \\
\hline
cc-pVDZ & 24 & $-76.2410$ & $+40.7$ & 107\,382 (46) & 674.3 \\
Slater DZP & 24 & $-76.2817$ & $0$ & \gvq{h2o_ccpvdz_pauli}{107\,382} (46) & \gvq{p62_dzp_lambda_fc}{706.3} \\
cc-pVDZ+TZpol & 35 & \gvq{p62_tzpol_fc}{$-76.2665$} & $+\gvq{p62_tzpol_above_dzp_mha}{15.2}$ & \gvq{p62_tzpol_pauli_fc}{505\,467} (68) & \gvq{p62_tzpol_lambda_fc}{1\,814.9} \\
aug-cc-pVDZ & 41 & \gvq{p62_aug_fc}{$-76.2737$} & $+\gvq{p62_aug_above_dzp_mha}{8.0}$ & \gvq{p62_aug_pauli_fc}{982\,785} (80) & \gvq{p62_aug_lambda_fc}{1\,901.1} \\
cc-pVTZ & 58 & $-76.3322$ & $-50.5$ & 3\,983\,652 (114) & 7\,653.5 \\
\end{tabular}
\end{ruledtabular}
\end{table}

\textbf{[MEASURED, pre-registered; GO]}  Table~\ref{tab:matched}.  Both
predictions held: the 35-function basis stays
$\gvq{p62_tzpol_above_dzp_mha}{15.2}\mHa$ above the Slater DZP and
aug-cc-pVDZ $\gvq{p62_aug_above_dzp_mha}{8.0}\mHa$ above it.  Both steps
are correlation-dominated, with mean-field shares of $0.19$ and $0.45$
against the Slater DZP step's $0.71$, and the two rungs keep neon's
order; those were the reported clauses, and they held too.  The
aug-cc-pVDZ energy agrees with the commonly quoted value near
$-76.274$~Ha, which is an orientation check and not a cited reference.

The crossing therefore lies between aug-cc-pVDZ and cc-pVTZ, $0.14$ of
the way in energy.  The lower bound needs no interpolation: aug-cc-pVDZ,
still $8$~mHa short of the Slater DZP, already costs
$\gvq{p62_aug_pauli_ratio}{9.2}\times$ the Pauli terms,
$\gvq{p62_aug_lambda_ratio}{2.7}\times$ the 1-norm, and 80 qubits
against 46.  Read linearly in energy between the bracketing rungs, the
matched-accuracy Gaussian cost is near $1.4$ million terms, a 1-norm
near $2\,700$~Ha and about 85 qubits, which is 13 times the terms and
3.8 times the 1-norm; that interpolation is a reading of two measured
rungs, not a measurement of any basis.  The qubit statement of this
paper is thus a matched-accuracy one: at matched or better accuracy on
water, the Gaussian basis costs at least nine times the Pauli terms and
2.7 times the 1-norm of the Slater DZP, and the equal-size result of
Table~\ref{tab:qubit} is its lower rung.

Two things the angular sparsity of the one-centre framework
\cite{paper22} would have added, and does not here.  Multi-centre
integrals over these bases are dense; no Gaunt zero appears in any
Hamiltonian of Table~\ref{tab:qubit}, on either side.  And the 1-norm
observation of Sec.~\ref{sec:native}, that a low 1-norm at equal size
tracks a worse energy, is the reason the composed library's 1-norm
figures (Papers~14, 20) are reported at matched qubits with the
accuracy stated alongside, never as a cost advantage on their own.

\section{Prior art}
\label{sec:prior}

A moderate search was made, against primary sources where reachable.
The mean-field half of Sec.~\ref{sec:dzp}, that Slater functions beat
Gaussians per function at the SCF level, is textbook and no part of it
is claimed here.  The design recipe, atom-optimized Slater primitives
generally contracted on the cc-pV$X$Z pattern, has precedent in the
Madrid group's STO valence sets for H to Ne with HF- and CI-optimized
exponents \cite{ema2003}, and in the correlation-consistent-style STO
sets of F\"orster and Visscher \cite{forster2021}.  The closest
molecular STO-against-cc-pV$X$Z comparison found compares forces, not
correlated energies \cite{rico2004}.  No matched-size correlated
STO-against-GTO head-to-head on a molecule was found, nor the split of
the DZ-to-TZ gap into its mean-field and correlated parts, nor any
qubit-resource comparison between the two families.  A qualitative
reading that orbital shape matters little for correlation energy is
sometimes attributed to Persson and Taylor; the full text could not be
read and it is not cited as support.  The water complete-basis anchor of
Helgaker, Klopper, Koch and Noga \cite{helgaker1997} was consulted for
orientation only; no basis-set-limit statement is made.

The transfer question of Sec.~\ref{sec:ladder}, whether exponents set on
the free atom carry to the molecule and at what ratio, is the standard
premise of every atom-optimized basis family, cc-pV$X$Z included
\cite{dunning1989}.  What is specific here is the pre-registered
rung-by-rung measurement of the ratio, the identification of the
$p$-split as the carrier on both sides, and the diagnosis of the one
rung where the ratio has the wrong sign.

\section{Limitations}
\label{sec:limits}

\begin{itemize}
\item One geometry, one method, one molecule.  No potential energy
surface, no basis-set limit, no second molecule.  The ladder frees
exponents in one order and does not test every partition of five
primitives.  ``Needed'' in Sec.~\ref{sec:ladder} means ``worth more
than the bound at twenty-four functions''; it does not mean a count is
forced.
\item The Slater functions are Gaussian fits.  The fit is converged at
the $0.07$~mHa level for the DZP basis and is exponent-independent by
construction on the ladder, but every number here is for the fitted
basis.  Exact Slater integrals on three centres exist in the series
(Paper~59) and were not used for these energies.
\item The matched-accuracy comparison (Sec.~\ref{sec:matched}) uses
standard Gaussian bases as published against the Slater DZP as built,
with its water-optimized polarization; cc-pVDZ's own re-optimized
polarization is worth $1.2$~mHa and does not change the bracket.  The
35-function intermediate basis is our construction.  The interpolated
matched-accuracy cost is a reading between two measured rungs.
\item The $3d$ exponent on the ladder is tied to the $3p$, and
Sec.~\ref{sec:r3} shows that tie to be the ladder's one design flaw.
A ladder with the $3d$ set by an atomic correlated criterion would be a
different ladder and has not been run.
\item The solver is the project's own.  It is validated against an
external reference on STO-3G water and against exact limits, but no
larger-basis external cross-check of the CCSD(T) energies exists on the
host.
\item None of this is a GeoVac lattice result, and none of it uses the
angular sparsity that is the framework's structural asset.  The paper
measures what tuned exponential functions buy at equal size, and finds
that the framework's fixed-exponent rules do not buy it.
\end{itemize}

\section{Conclusion}
\label{sec:conclusion}

At twenty-four functions on water, a Slater-type basis built by
cc-pVDZ's own rule is $41$~mHa lower at frozen-core CCSD(T), mostly at
the mean-field level, for the same Pauli count and five percent more
1-norm.  The gain is exponent choice: the one-centre framework's
fixed-exponent rules, hydrogenic $Z/n$ and single-scale Sturmian, lose
by $0.2$ to more than $1$~Ha at the same size, and the hydrogenic loss is
the node structure the fixed exponents impose, which screening makes
worse.  Releasing the oxygen exponents one at a time, each set on the
free atom, closes that gap: five exponents put water $3$~mHa below
cc-pVDZ with nothing tuned on the molecule, the $p$ split carries nine
tenths of the gain on both the atom and the molecule, and the gain
transfers at about half, all as predicted before the energies were
computed.  The one rung that went the wrong way is a $3d$ passenger the
atomic objective cannot see, plus two mean-field pieces that cancel on
water and not on the atom.  At matched accuracy the comparison turns
around: aug-cc-pVDZ at 41 functions is still $8$~mHa short of the
Slater DZP and already costs nine times the Pauli terms and 2.7 times
the 1-norm, and the crossing lies between it and cc-pVTZ.  The
water-tuned basis stays $34$~mHa ahead of the best atom-set rung, and
nothing here is a property of the lattice.

\section*{Backing}

Solver, validation tests, the radial atom, the kit builders and the
Pauli counter: \texttt{tests/test\_water\_ccsd.py}.  The ladder's guard,
\texttt{tests/test\_paper62\_scale\_ladder.py}, recomputes the
free-atom objective at the stored exponents and checks that each rung's
exponents are a stationary minimum of it, derives the SL1 and SL2
ratios from the stored water energies and the recomputed atom energies,
and recomputes the one-centre $R3$ probe of Sec.~\ref{sec:r3}; only
the water energies are read from its data file.  The matched-accuracy
guard, \texttt{tests/test\_paper62\_matched\_accuracy.py}, recomputes
the aug-cc-pVDZ and cc-pVDZ water runs end to end, energies and
Jordan--Wigner cost, from basis files stored beside it; the Slater DZP
target and the other two rungs are read as registered values.  The ladder's record,
with its predictions, disposition and the $R3$ addendum:
\texttt{docs/predictions/2026-10-30\_scale\_ladder.md}; the
matched-accuracy record:
\texttt{docs/predictions/2026-10-31\_matched\_accuracy\_gaussian.md}.
Chronicle: CHANGELOG v5.20.11, v5.20.12, v5.21.25, v5.21.28, v5.21.29.

\begin{thebibliography}{99}

\bibitem{avery2006}
J.~E.~Avery and J.~S.~Avery, \textit{Generalized Sturmians and Atomic Spectra}
(World Scientific, Singapore, 2006).

\bibitem{avery2017}
J.~E.~Avery and J.~S.~Avery, \textit{Hyperspherical Harmonics and Their
Physical Applications} (World Scientific, Singapore, 2017),
doi:10.1142/10690.

\bibitem{dunning1989}
T.~H.~Dunning, Jr., ``Gaussian basis sets for use in correlated molecular
calculations. I. The atoms boron through neon and hydrogen,''
J.\ Chem.\ Phys.\ \textbf{90}, 1007 (1989).

\bibitem{mcmurchie1978}
L.~E.~McMurchie and E.~R.~Davidson, ``One- and two-electron integrals over
Cartesian Gaussian functions,'' J.\ Comput.\ Phys.\ \textbf{26}, 218 (1978).

\bibitem{stewart1970}
R.~F.~Stewart, ``Small Gaussian expansions of Slater-type orbitals,''
J.\ Chem.\ Phys.\ \textbf{52}, 431 (1970).

\bibitem{purvis1982}
G.~D.~Purvis III and R.~J.~Bartlett, ``A full coupled-cluster singles and
doubles model: The inclusion of disconnected triples,''
J.\ Chem.\ Phys.\ \textbf{76}, 1910 (1982).

\bibitem{raghavachari1989}
K.~Raghavachari, G.~W.~Trucks, J.~A.~Pople, and M.~Head-Gordon, ``A fifth-order
perturbation comparison of electron correlation theories,''
Chem.\ Phys.\ Lett.\ \textbf{157}, 479 (1989).

\bibitem{crawford_projects}
T.~D.~Crawford and co-workers, \textit{Programming Projects}, the
STO-3G water CCSD(T) reference values, \url{https://github.com/CrawfordGroup/ProgrammingProjects}.
[Cited for the reference energies the solver is validated against.]

\bibitem{slater1960}
J.~C.~Slater, \textit{Quantum Theory of Atomic Structure}, Vol.~I
(McGraw-Hill, New York, 1960).
[Cited for the $F^{k}$ and $G^{k}$ Slater integrals and the
average-of-configuration energy.]

\bibitem{jordan1928}
P.~Jordan and E.~Wigner, ``\"Uber das Paulische \"Aquivalenzverbot,''
Z.\ Phys.\ \textbf{47}, 631 (1928).

\bibitem{trenev2025}
D.~Trenev, P.~J.~Ollitrault, S.~M.~Harwood, T.~P.~Gujarati,
S.~Raman, A.~Mezzacapo, and S.~Mostame,
``Refining resource estimation for the quantum computation of
vibrational molecular spectra through Trotter error analysis,''
Quantum \textbf{9}, 1630 (2025); arXiv:2311.03719.
[Cited for the frozen-core cc-pVDZ water Pauli count of Table~5,
which the counter here reproduces exactly.]

\bibitem{ema2003}
I.~Ema, J.~M.~Garc\'ia de la Vega, G.~Ram\'irez, R.~L\'opez,
J.~Fern\'andez Rico, H.~Meissner, and J.~Paldus, ``Polarized basis sets of
Slater-type orbitals: H to Ne atoms,'' J.\ Comput.\ Chem.\ \textbf{24},
859 (2003), doi:10.1002/jcc.10227.

\bibitem{forster2021}
A.~F\"orster and L.~Visscher, ``GW100: A Slater-type orbital perspective,''
J.\ Chem.\ Theory Comput.\ \textbf{17}, 5080 (2021),
doi:10.1021/acs.jctc.1c00308.

\bibitem{rico2004}
J.~Fern\'andez Rico, R.~L\'opez, I.~Ema, and G.~Ram\'irez, ``Reference
calculations with Slater-type orbitals: forces,'' Int.\ J.\ Quantum
Chem.\ \textbf{100}, 221 (2004), doi:10.1002/qua.20088.

\bibitem{helgaker1997}
T.~Helgaker, W.~Klopper, H.~Koch, and J.~Noga, ``Basis-set convergence of
correlated calculations on water,'' J.\ Chem.\ Phys.\ \textbf{106}, 9639
(1997), doi:10.1063/1.473863.

\bibitem{paper22}
J.~Loutey, ``Potential-Independent Angular Sparsity for Qubit Hamiltonians in Spherical Fermion Systems'' (Paper 22 in this series).

\bibitem{sl_record}
J.~Loutey, pre-registered record ``How many radial scales water needs on
oxygen (SL),'' \texttt{docs/predictions/2026-10-30\_scale\_ladder.md}, this
repository, registered and committed 2026-10-30 before the scored
energies were computed; disposition and the $R3$ addendum appended the
same day.

\bibitem{ma_record}
J.~Loutey, pre-registered record ``The matched-accuracy Gaussian
comparison for water (MA),''
\texttt{docs/predictions/2026-10-31\_matched\_accuracy\_gaussian.md},
this repository, registered and committed 2026-10-31 before the two
scored energies were computed; disposition appended the same day.

\end{thebibliography}

\end{document}
